\ifdefined\HANDOUTMODE
  \PassOptionsToClass{handout}{beamer}
\fi
\documentclass[10pt,aspectratio=169]{beamer}
\usetheme{Warsaw}

\title{Pallas Storage Compression Report}
\author{Yash Rawat}
\institute{Telecom SudParis, France}
\date{\today}

\begin{document}

\begin{frame}
  \titlepage
\end{frame}

\begin{frame}{Overview}
  \begin{itemize}
    \item<1-> Motivation: reduce the storage cost of redundant timing data in Pallas traces without compromising the usefulness of the recorded execution behavior.
    \item<2-> Trace shaping: improved loop and sequence detection so repeated polling patterns can be represented more compactly.
    \item<3-> Lossless encoding: evaluated timestamp and duration compression schemes that preserve exact values while reducing trace size.
    \item<4-> Lossy encoding: introduced controlled approximation paths for highly redundant timing vectors to study compression-quality tradeoffs.
    \item<5-> Benchmarking: built write-time, read-time, verification, and post-mortem analysis tooling to quantify overheads and compression gains.
  \end{itemize}
\end{frame}

\begin{frame}[fragile]{Motivation}
  \begin{columns}[T]
    \begin{column}{0.5\textwidth}
      \begin{block}{Spinloop Pattern}
\small
\begin{verbatim}
for iter = 1 .. iterations:
  MPI_Barrier(comm)

  if rank == receiver:
    MPI_Irecv(..., &req)
    flag = 0
    while !flag:
      MPI_Test(&req, &flag, IGNORE)
  else:
    sleep(delay_us)
    MPI_Isend(..., &req)
    MPI_Wait(&req, IGNORE)
\end{verbatim}
      \end{block}
    \end{column}
    \begin{column}{0.48\textwidth}
      \begin{itemize}
        \item<2-> Each delayed send forces the receiver into a hot polling loop with many failed \texttt{MPI\_Test} calls before one final success.
        \item<3-> Across large iteration counts, this generates massive amounts of repetitive timing data with very similar sequence structure.
        \item<4-> The resulting Pallas traces are often dominated by timestamp and duration vectors attached to these repeated polling sequences.
        \item<5-> Our goal is to compress this redundancy while preserving the execution information needed for post-mortem analysis.
      \end{itemize}
    \end{column}
  \end{columns}
\end{frame}

\begin{frame}{Supporting Data}
  \small
  \begin{block}{Spinloop Benchmark Cases}
    \centering
    \begin{tabular}{|c|c|c|}
      \hline
      Iterations & Delay ($\mu$s) & Trace Directory Size \\
      \hline
      1000 & 10000 & \textit{fill} \\
      \hline
      10000 & 10000 & \textit{fill} \\
      \hline
      10000 & 10000 & \textit{fill} \\
      \hline
      10000 & 100000 & \textit{fill} \\
      \hline
    \end{tabular}
  \end{block}

  \vspace{0.5em}
  \begin{itemize}
    \item Use this table to connect the workload parameters directly to the growth of the generated trace directory.
    \item The directory size can later be compared against compression ratio, read overhead, and write overhead plots.
  \end{itemize}
\end{frame}

\end{document}
